\documentclass[10pt, a4paper]{article}

% --- paquets et configuration ---
\usepackage[utf8]{inputenc}
\usepackage[T1]{fontenc}
\usepackage[french]{babel}
\usepackage[margin=1.8cm]{geometry}
\renewcommand{\familydefault}{\sfdefault}
\usepackage{xcolor}
\usepackage{titlesec}
\usepackage{enumitem}
\usepackage{tcolorbox}
\tcbuselibrary{skins}

% --- couleurs ---
\definecolor{desczone}{RGB}{45, 90, 120}
\definecolor{desccreature}{RGB}{120, 50, 70}
\definecolor{descobstacle}{RGB}{140, 100, 40}
\definecolor{descpuzzle}{RGB}{60, 110, 80}
\definecolor{desctransition}{RGB}{80, 70, 110}
\definecolor{descmoment}{RGB}{100, 80, 60}

% configuration des titres
\titleformat{\section}{\large\bfseries\color{desczone}\uppercase}{}{0em}{}[\hrule]
\titleformat{\subsection}{\normalsize\bfseries\color{desccreature}}{}{0em}{}
\titleformat{\subsubsection}{\small\bfseries}{}{0em}{}

% boîtes de description
\newtcolorbox{zonebox}[1]{colback=desczone!5, colframe=desczone, title=\textbf{#1}, sharp corners, fonttitle=\small}
\newtcolorbox{creaturebox}[1]{colback=desccreature!5, colframe=desccreature, title=\textbf{#1}, sharp corners, fonttitle=\small}
\newtcolorbox{obstaclebox}[1]{colback=descobstacle!5, colframe=descobstacle, title=\textbf{#1}, sharp corners, fonttitle=\small}
\newtcolorbox{puzzlebox}[1]{colback=descpuzzle!5, colframe=descpuzzle, title=\textbf{#1}, sharp corners, fonttitle=\small}
\newtcolorbox{transitionbox}[1]{colback=desctransition!5, colframe=desctransition, title=\textbf{#1}, sharp corners, fonttitle=\small}
\newtcolorbox{momentbox}[1]{colback=descmoment!5, colframe=descmoment, title=\textbf{#1}, sharp corners, fonttitle=\small}

\setlist[itemize]{label=\textbullet, leftmargin=*, noitemsep, topsep=2pt}

\begin{document}

\vspace{0.5cm}

% ============================================================
\section*{jour 1 -- le marais}
% ============================================================

\begin{zonebox}{arrivée -- le réveil dans la capsule}
\textit{« Vous sentez quelque chose entre le soufre, le métal chauffé, et le végétal qui pourrit. L'odeur vous frappe avant la lumière. À travers le hublot fissuré, des formes bougent dans la brume. Certaines luisent. D'autres semblent respirer. La capsule craque. Votre casque est lourd et vous avez l'impression d'étouffer. Est-que que c'est respirable ici ? »}
\end{zonebox}

\begin{zonebox}{ambiance -- le marais}
\textit{« Sol qui cède sous les pas avec un bruit de succion. Air épais, saturé de spores. Bulles qui remontent des flaques. Cris au loin, impossibles à identifier. »}
\end{zonebox}

\begin{creaturebox}{première apparition -- le patriarche-composteur}
\textit{« Ce que vous preniez pour une colline se lève. Quatre étages de bois humide, de mousse et de champignons grouillants. Ses doigts -- des troncs entiers -- se referment sur la capsule et la soulèvent comme on cueillerait une fleur. Puis il s'éloigne, et la capsule avec. Il la dépose devant l'entrée d'une grotte qui commence à la décortiquer en morceaux. »}
\end{creaturebox}

\begin{obstaclebox}{découverte -- la grotte aux cadavres}
\textit{« Vous vous approchez de la grotte. L'air qui en sort est âcre. Vos yeux s'habituent à la pénombre et vous les voyez : des corps humains. Des excroissances fongiques ont poussé sur eux, délicates comme des fleurs, luminescentes comme des bijoux. De petits écosystèmes prospèrent sur chacun d'eux. »}
\end{obstaclebox}

\begin{puzzlebox}{puzzle -- la traversée de la grotte}
\textit{« Des vessies translucides pulsent aux parois, gonflant et dégonflant lentement, expulsant à chaque expiration un gaz bleuté. L'air a un goût sucré. De l'autre côté du gouffre, la sortie -- une pente qui remonte vers la lumière. Entre vous et elle : vingt mètres de brume et quelque chose de massif, immobile dans l'ombre. »}

\textit{(on le dit si jet de perception correct)} \textit{« Les vessies semblent filtrer quelque chose. Leur membrane est d'un violet profond, et rien ne pousse directement autour d'elles. »}
\end{puzzlebox}

\begin{momentbox}{découverte -- l'outre-pulmonaire}
\textit{« Tiède dans vos mains. La membrane vibre, émet une faible lumière violette. L'air qui en sort, au contact de votre peau, semble frais. La vessie pulse contre vous, comme s'y elle cherchait à s'attacher. »}
\end{momentbox}

\begin{momentbox}{découverte -- l'anguille-câble}
\textit{« Dans l'eau sombre, un ver tubulaire translucide d'environ un mètre. Son système nerveux dessine des arabesques dorées sous la peau. Il est emmêlé dans des racines, frémit faiblement. Quand votre main s'approche, un picotement, et vous la retirez aussitôt, par réflexe, nerveusement. »}
\end{momentbox}

\begin{creaturebox}{première apparition -- le crapaud-hallucinogène}
\textit{« Vous ne l'avez pas vu bouger. Il était là depuis le début. Une masse bulbeuse tapie dans un recoin, sa peau secrète un film irisé qui s'évapore en volutes paresseuses. »}
\end{creaturebox}

\begin{momentbox}{effet -- les hallucinations commencent}
\textit{(choisir selon le personnage)}
\begin{itemize}
    \item \textit{« Tu baisses les yeux sur tes mains. Sous ta peau, quelque chose bouge. Des vers, peut-être. »}
    \item \textit{« Le cadavre à ta gauche vient de cligner des yeux. Non. Il te sourit. Tu connais ce sourire. »}
    \item \textit{« Le sol devient mou. Tu t'enfonces. Personne ne semble remarquer. »}
    \item \textit{« Ta propre voix crie ton nom. Ta bouche est fermée. »}
    \item \textit{« Un de tes compagnons a le visage de ta petite sœur. Tu n'as pas de petite sœur. »}
\end{itemize}
\end{momentbox}

\begin{obstaclebox}{découverte -- l'étendue aux méduses}
\textit{« La surface de l'eau est parsemée de dômes translucides -- des méduses, par centaines. La brume acide qu'elles sécrètent fait onduler l'air. De l'autre côté, un sol plus ferme. »}
\end{obstaclebox}

\begin{puzzlebox}{puzzle -- les bisons-creusets}
\textit{« Un grondement sourd. Des silhouettes massives émergent de la brume -- quadrupèdes de la taille d'éléphants, peau pareille à de la croûte volcanique craquelée. Des fentes sur leurs flancs luisent d'orange, comme des braises sous la cendre. De petites créatures glabres entrent et sortent de ces ouvertures, déposent des sphères luisantes au sol, y retournent. Les bêtes avancent dans l'eau. Les méduses s'écartent sur leur passage. Un seul bison est resté en retrait de votre côté du marais. On ne saurait pas bien dire ce qu'il est en train de faire. »}
\end{puzzlebox}

\begin{creaturebox}{première apparition -- le bison-creuset}
\textit{« Vue de près, une centrale thermique vivante. Sa respiration fait vibrer l'air. Ses yeux -- deux billes de charbon incandescent -- vous traversent sans vous voir. »}
\end{creaturebox}

\begin{momentbox}{découverte -- la sphère galvanique}
\textit{« Un des petits a déposé sa charge au sol avant de retourner dans les flancs de sa mère. Une sphère grosse comme un pamplemousse, orange pulsante. Même à un mètre, la chaleur. L'air autour ondule. »}
\end{momentbox}

\begin{transitionbox}{nuit 1 -- installation du camp}
\textit{« La température chute. Les champignons au loin commencent à luire d'un bleu irréel. Le dôme vous isole du monde, pas de ses sons : gargouillis, craquements, et ces pas lourds qui font trembler le sol. »}
\end{transitionbox}

\begin{momentbox}{événement nocturne -- les pagures géants}
\textit{« Le sol vibre en rythme, comme un cœur de pierre. Par la paroi translucide du dôme, des rochers qui marchent. Une file de créatures massives -- chitine et calcaire -- avançant avec l'implacable patience des ères géologiques. Leur chemin passe exactement là où vous avez planté votre camp. »}
\end{momentbox}

\begin{creaturebox}{première apparition -- le pagure géant}
\textit{« Un rocher de la taille d'une voiture. Avec des pattes -- une douzaine, trapues comme des piliers. Il avance. Lentement. »}
\end{creaturebox}

% ============================================================
\section*{jour 2 -- la forêt fongique}
% ============================================================

\begin{zonebox}{arrivée -- la forêt fongique}
\textit{« Une canopée de champignons géants filtre une lumière bleutée, douce. Le sol spongieux s'illumine sous vos pas -- mousse phosphorescente qui dessine vos empreintes. L'air est frais, parfumé. Des chants d'oiseaux-insectes dans le silence feutré. Pas un bruit de prédateur. »}
\end{zonebox}

\begin{obstaclebox}{découverte -- la muraille végétale}
\textit{« Le chemin s'arrête net. Dix mètres de ronces noires, denses comme du béton, à perte de vue dans les deux directions. Les épines luisent d'une substance humide. »}
\end{obstaclebox}

\begin{puzzlebox}{puzzle -- l'écaille-clé}
\textit{« Plus loin, un mastodonte herbivore couvert d'écailles osseuses traverse la forêt. Devant lui, la muraille s'ouvre -- les ronces s'écartent comme des rideaux et se referment derrière. En passant, il se frotte contre un rocher. Quelques écailles tombent. Les ronces autour frémissent. »}
\end{puzzlebox}

\begin{creaturebox}{première apparition -- le bison-cuirassé}
\textit{« Il mange. Sa mâchoire broie la végétation. Ses yeux vous trouvent un instant -- deux billes ternes -- puis vous oublient. »}
\end{creaturebox}

\begin{creaturebox}{première apparition -- le gardien-lémurien}
\textit{« Deux yeux immenses au-dessus de vous -- yeux de lémurien, curieux et anciens -- depuis une petite branche. Cinquante centimètres à peine. Des veines vertes luminescentes courent sous sa peau et semblent se prolonger dans la muraille. Quelque chose brille dans sa main : un de vos boutons, peut-être. Elle le fait tourner entre ses doigts. »}
\end{creaturebox}

\begin{momentbox}{jeu de mime -- le gardien-lémurien}
\textit{« La créature descend avec une agilité fluide. Sans prévenir : trois tapotements sur son torse, un mouvement circulaire de la main, un saut sur place. Puis immobile. Elle vous fixe. »}
\end{momentbox}

\begin{obstaclebox}{découverte -- la cabane d'oren}
\textit{« Entre deux champignons géants, une cabane dans le tronc évidé d'un champignon mort. Un potager luxuriant -- légumes impossibles, fruits trop parfaits, couleurs trop vives. De la fumée sort d'une cheminée. »}
\end{obstaclebox}

\begin{creaturebox}{première apparition -- oren le parasité}
\textit{« La porte s'ouvre sur un visage humain. Presque. La peau a une teinte verdâtre, subtile. "Des visiteurs ! Ça faisait tellement longtemps." Sa voix est douce. Quand il bouge, sous sa peau, des filaments ondulent à la lumière. »}
\end{creaturebox}

\begin{momentbox}{indices de la condition d'oren}
\textit{(à distiller au fil de la conversation)}
\begin{itemize}
    \item \textit{« Quand il tend la main vers la lumière, des filaments verts pulsent sous sa peau. »}
    \item \textit{« Il parle du vaisseau comme d'un souvenir lointain. "Le vide froid de l'espace... pourquoi y retourner ?" »}
    \item \textit{« Le robot émet un bourdonnement quand Oren s'approche. Ses capteurs semblent perturbés. »}
    \item \textit{« Ses plantes se tournent vers lui quand il passe. »}
\end{itemize}
\end{momentbox}

\begin{creaturebox}{découverte -- l'écureuil-coffre}
\textit{« Un petit rongeur sort des fougères. Fourrure élastique, presque caoutchouteuse. Grands yeux curieux. Il renifle l'air, mâchouille des graines. Il ne fuit pas. »}
\end{creaturebox}

\begin{momentbox}{l'écureuil-coffre gonflé}
\textit{« L'écureuil gonfle. Sa peau s'étire, révèle une membrane translucide sous la fourrure. Ses pattes quittent le sol. Il flotte, rebondit mollement contre les branches, vous fixe de ses grands yeux. »}
\end{momentbox}

\begin{momentbox}{phrases d'oren -- hospitalité}
\textit{(à utiliser en roleplay)}
\begin{itemize}
    \item \textit{« Vous êtes les premiers visiteurs depuis... combien de temps déjà ? Le temps n'a plus d'importance ici. »}
    \item \textit{« Le vaisseau ? Ah oui... Pourquoi voudriez-vous retourner dans le vide froid de l'espace ? »}
    \item \textit{« Goûtez ces fruits. La planète nous donne tout ce dont nous avons besoin. »}
    \item \textit{« Vous semblez fatigués. Restez cette nuit. »}
    \item \textit{« Je ne suis pas seul ici. La forêt prend soin de moi. »}
\end{itemize}
\end{momentbox}

\begin{momentbox}{combat -- oren se transforme}
\textit{« Ses yeux noircissent. Lentement, comme de l'encre dans de l'eau. La mousse sous sa peau éclate à la surface -- armure végétale qui craque et grince. Des lianes jaillissent de ses paumes, de son dos, du sol. Sa voix résonne comme mille voix superposées. "Je voulais vous aider." »}
\end{momentbox}

\begin{obstaclebox}{découverte -- la rivière de sève acide}
\textit{« Un cours d'eau qui n'en est pas un. Liquide ambré, épais, qui dissout lentement tout ce qu'il touche. Ossements blanchis sur les berges. Racines à moitié digérées. Cinq mètres de large. Le chemin continue de l'autre côté. »}
\end{obstaclebox}

\begin{puzzlebox}{puzzle -- les libélucioles}
\textit{« Il y a de larges champignons qui flottent au-dessus de la rivière. On voit aussi des insectes volants. Exosquelette plein de prises naturelles, ailes vrombissantes. Seuls, trop légers pour porter quoi que ce soit. Mais ils s'accrochent parfois les uns aux autres, forment des grappes flottantes. Un point surélevé domine la rivière de l'autre côté. »}
\end{puzzlebox}

\begin{momentbox}{traversée -- avec les libélucioles}
\textit{« Le saut. La chute. Puis les insectes -- une grappe vrombissante accrochée à vos vêtements, vos cheveux, vos bras. Vous planez, maladroitement, au-dessus du liquide ambré. L'autre rive approche. Les insectes fatiguent. Vous perdez de l'altitude. Le sol. »}
\end{momentbox}

\begin{momentbox}{échec -- éclaboussure de sève}
\textit{« La sève touche votre bras. La peau fume, se dissout couche par couche. La douleur est blanche. »}
\end{momentbox}

\begin{creaturebox}{première apparition -- le tisseur chirurgical}
\textit{« Un cliquetis rapide, métallique, comme une machine à coudre. Une araignée de mer décharnée, haute sur huit pattes-aiguilles d'os blanc. Son abdomen translucide pulse d'un liquide soyeux. Elle s'arrête, tâte l'air. Se tourne vers celui d'entre vous qui saigne. »}
\end{creaturebox}

\begin{momentbox}{le tisseur soigne -- description}
\textit{« Ses pattes-aiguilles perforent la chair autour de la plaie. La douleur est immédiate. Point après point, un fil organique referme la blessure. Proprement. »}
\end{momentbox}

\begin{transitionbox}{nuit 2 -- la forêt s'illumine}
\textit{« La nuit tombe. Des millions de spores dorées flottent dans l'air. Les champignons pulsent d'un bleu profond. »}
\end{transitionbox}

\begin{momentbox}{événement nocturne -- la ruche-anémone}
\textit{« Une lueur. Douce. Elle pulse dans l'obscurité. Vos jambes vous portent vers elle avant que vous ne décidiez de vous lever. Anémone géante, tentacules lumineux ondulant lentement. Vos mains se tendent. Les tentacules s'enroulent autour de vos doigts, doux comme de la soie. Tiède. »}
\end{momentbox}

\begin{momentbox}{réveil forcé -- la ruche-anémone}
\textit{(si un autre joueur intervient)}
\textit{« Le coup vous réveille. Vos mains, à moitié englouties par des tentacules lumineux. Votre peau couverte de marques de succion. »}
\end{momentbox}

% ============================================================
\section*{jour 3 -- le rivage de mer}
% ============================================================

\begin{zonebox}{arrivée -- le rivage et l'orage}
\textit{« Sable gris. Ciel de cuivre. L'air chargé d'ozone fait dresser les cheveux sur vos nuques. Par-delà le bras de mer, les montagnes. Derrière elles, le signal du GPS : le vaisseau. La plage est par endroits jonchée d'oeufs rosâtres et diaphanes qui ont la teinte douce des nouveaux-nés en bonne santé. »}
\end{zonebox}

\begin{momentbox}{annonce -- l'orage de silice}
\textit{« Le robot s'immobilise. "Alerte météorologique critique. Orage de silice. Impact : trente minutes. Recommandation : abri immédiat." Sur l'horizon, une muraille de nuages noirs striés d'éclairs horizontaux. À l'intérieur, des rochers, des arbres et des éclats de cristal tournoient. »}
\end{momentbox}

\begin{obstaclebox}{découverte -- le bras de mer}
\textit{« L'eau s'étend, noire et calme. Impossible de voir le fond. La rive opposée : cinq cents mètres. »}
\end{obstaclebox}

\begin{puzzlebox}{puzzle -- les baleines-méduses}
\textit{« Des silhouettes émergent près du rivage. Colossales, translucides -- entre la baleine et la méduse, des cathédrales de gelée qui pulsent doucement. Des dizaines, gros comme des pastèques. Puis l'orage approche. La créature déploie une trompe massive et commence à aspirer les œufs que vous voyiez plus tôt, un par un, et les emporte dans son ventre. Un lézard s'est rué sur un des oeufs sur le point d'être gobé comme un flamby »}
\end{puzzlebox}

\begin{creaturebox}{la baleine-méduse si on veut la décrire}
\textit{« Sa surface ondule comme de la soie liquide, veines bioluminescentes pulsant sous la gelée. Elle glisse devant vous sans ralentir ni s'écarter. »}
\end{creaturebox}

\begin{momentbox}{l'aspiration}
\textit{« La trompe vous trouve. Vous êtes aspiré dans un tunnel de chair douce, propulsé à travers des membranes qui cèdent comme des rideaux. Une cavité baignée de lumière bleue. Des dizaines d'œufs flottent autour de vous. L'air est respirable -- la trompe fait office de tuba géant. À travers les parois, l'océan défile : créatures étranges, bancs de poissons étranges, silhouettes massives dans les profondeurs étranges. »}
\end{momentbox}

\begin{creaturebox}{première apparition -- le crabe-foreur-cafard}
\textit{« Un bruit de grattement dans la trompe. Mi-crabe, mi-cafard, mandibules foreuses tournant à vide, queue de scorpion dressée. »}
\end{creaturebox}

\begin{momentbox}{combat dans la baleine -- ambiance}
\textit{« Le sol ondule sous vos pieds. L'espace est confiné, les œufs roulent et rebondissent. La bile acide du crabe creuse des sillons fumants dans la chair de la baleine. »}
\end{momentbox}

\begin{momentbox}{l'arrivée sur l'autre rive}
\textit{« Une contraction vous projette à travers les membranes. Sable gris. Trempés, couverts de mucus. Derrière, la tempête -- éclairs horizontaux, vagues de cristaux. Devant, les montagnes. »}
\end{momentbox}

\begin{transitionbox}{nuit 3 -- dans la tempête}
\textit{« La tempête hurle toute la nuit. Dans les flashs d'éclairs, des silhouettes -- des créatures qui volent dans l'ouragan, surfent sur les vents. Le robot ne dort pas. Ses capteurs fixent le ciel. »}
\end{transitionbox}

% ============================================================
\section*{jour 4 -- la montagne aride}
% ============================================================

\begin{zonebox}{arrivée -- la montagne aride}
\textit{« L'air sec comme un four. Le terrain monte, escarpé, ocre et gris. Vos lèvres se craquellent. Et là-haut, au sommet du plateau, le vaisseau. Un point métallique dans la lumière crue. »}
\end{zonebox}

\begin{momentbox}{la soif -- progression}
\textit{(à lire selon le temps écoulé sans eau)}
\begin{itemize}
    \item \textit{5 min : « Bouche sèche. »}
    \item \textit{10 min : « Votre langue colle à votre palais. Chaque respiration râpe. »}
    \item \textit{15 min : « Vos pensées deviennent floues. Le vaisseau danse dans la chaleur. »}
    \item \textit{20 min : « Vos muscles crient. Votre tête pulse. »}
\end{itemize}
\end{momentbox}

\begin{obstaclebox}{découverte -- l'oasis-tuyau}
\textit{« Au détour d'un rocher : un arbre blanc, cylindrique comme un tuyau géant planté dans le sol. À l'intérieur, de l'eau. Une source cristalline. Des fruits de toutes les couleurs. Des créatures de toutes tailles côte à côte -- proies et prédateurs, buvant à la même source. »}
\end{obstaclebox}

\begin{puzzlebox}{puzzle -- les règles de l'oasis}
\textit{« Les créatures s'approchent, boivent une gorgée, cueillent un fruit, puis s'écartent. Jamais deux gorgées. Jamais deux fruits. Au centre de l'oasis, une coupole chitineuse rouge-orangée pulse sous le couvert des feuilles. Immobile. »}
\end{puzzlebox}

\begin{creaturebox}{première apparition -- la coccisabre (piège)}
\textit{« Coupole chitineuse rouge-orangée, surmontée de boules de poils blancs en suspension statique. Elle n'a pas bougé depuis que vous l'observez. »}
\end{creaturebox}

\begin{creaturebox}{la coccisabre s'éveille}
\textit{« La coupole s'ouvre d'un coup. À l'intérieur, une lame d'un blanc d'os, longue d'un mètre. Un sifflement fait vibrer l'air. Les boules duveteuses dérivent vers vous. »}
\end{creaturebox}

\begin{obstaclebox}{découverte -- la falaise aux mille lames}
\textit{« La falaise : des arêtes cristallines au lieu de roches, tranchantes comme des rasoirs, luisant d'un éclat minéral. Le vaisseau est là-haut. »}
\end{obstaclebox}

\begin{puzzlebox}{puzzle -- les vers-colle}
\textit{« Sur la paroi, des vers translucides, épais comme des bras, grimpent lentement. Leur peau laisse une traînée argentée -- substance épaisse, visqueuse. Là où ils passent, les arêtes ne coupent plus. »}

\textit{(avertissement subtil)} \textit{« Un insecte se pose sur la traînée de mucus. Ses pattes collent. Il se débat. Il reste coincé. »}
\end{puzzlebox}

\begin{creaturebox}{première apparition -- les vers-colle}
\textit{« Translucides, doux au toucher. Si vous les caressez, ils frémissent et sécrètent une perle de mucus argenté. »}
\end{creaturebox}

\begin{momentbox}{le mucus -- conséquences}
\textit{« La substance argentée ne part pas. Ce que vous touchez colle. Votre arme refuse de quitter votre paume. Votre sac s'est transformé en boule compacte. »}
\end{momentbox}

\begin{creaturebox}{première apparition -- la pie-llard}
\textit{« Bourdonnement électrique au-dessus de vous. Plumage de débris métalliques cliquetant comme des pièces de monnaie. Une glande bleutée pulse sur son poitrail, crachotant des arcs d'électricité statique. Ses yeux se posent sur vos affaires. »}
\end{creaturebox}

\begin{obstaclebox}{découverte -- les arbres-cadavres}
\textit{« Des structures osseuses dressées comme des arbres. À leur sommet, des fruits orangés. Un fourmilier à deux trompes en engloutit un. Quelques secondes plus tard, des bipèdes aux queues chromées l'encerclent et l'emportent vers les hauteurs à toute vitesse. »}
\end{obstaclebox}

\begin{creaturebox}{première apparition -- les vélociraptors solaires}
\textit{« Ils chassent en meute, queues en éventail déployées vers le soleil. Ces queues -- écailles chromées comme des miroirs -- concentrent la lumière. Des marques de brûlure sur les rochers. Des arbres calcinés. »}
\end{creaturebox}

\begin{momentbox}{attaque -- convergence solaire}
\textit{« Un cercle parfait autour de vous. Les queues s'orientent. Des faisceaux de lumière convergent vers l'endroit où vous vous tenez. Le sol commence à fumer. »}
\end{momentbox}

\begin{momentbox}{tactique -- le fruit maudit}
\textit{« Le fruit explose. Une odeur si puissante qu'elle donne le vertige. Les vélociraptors se figent. Narines frémissantes. Puis, comme un seul être, ils se tournent vers la source et s'élancent. »}
\end{momentbox}

\begin{momentbox}{tactique -- se laisser capturer}
\textit{« Les serres vous saisissent. Le sol disparaît. Le vent hurle. Le vélociraptor vous emporte vers les hauteurs. Le nid est au sommet. Le sommet est là où vous devez aller. »}
\end{momentbox}

\begin{transitionbox}{nuit 4 -- le robot contemple}
\textit{« Le vaisseau est à portée, sa silhouette contre les étoiles. Mais le robot fixe le ciel, capteurs tournés ailleurs. "Pourquoi les humains regardent-ils les étoiles ?" Un silence. "Je comprends maintenant." »}
\end{transitionbox}

\begin{creaturebox}{événement nocturne -- l'écho-écho traque}
\textit{« Un silence. Pas le silence normal -- un silence actif. Puis un claquement de dents. Rapide. Proche. Une silhouette vibre dans la pénombre. Humanoïde. Pas de visage -- une parabole organique, une antenne de chair, tournée vers vous. »}
\end{creaturebox}

\begin{momentbox}{jeu du silence -- règles}
\textit{(annoncer aux joueurs)}
\textit{« Cette créature chasse au son. Elle est aveugle, mais elle entend tout. À partir de maintenant, vous devez chuchoter. En vrai. Si quelqu'un parle à voix normale, son personnage a fait du bruit. Et elle viendra. »}
\end{momentbox}

\begin{momentbox}{l'écho-écho attaque}
\textit{« Elle était là-bas. Maintenant elle est sur vous. Son cri -- une lame de son pur. Vos oreilles saignent. Votre équilibre s'effondre. »}
\end{momentbox}

% ============================================================
\section*{jour 5 -- la prairie et le vaisseau}
% ============================================================

\begin{zonebox}{arrivée -- la prairie}
\textit{« L'herbe haute ondule sous le vent. Des fleurs géantes ponctuent le paysage. Au-delà, vous voyez le vaisseau. Quoique pas très bien, parce que les herbes sont vraiment très hautes, et il est difficile d'imaginer les franchir. »}
\end{zonebox}

\begin{obstaclebox}{découverte -- la migration}
\textit{« Le sol tremble. Des milliers de quadrupèdes rayés, bulbes dorsaux pulsant en rythme, avancent en marée vivante. Ils couvrent tout l'horizon et déracinent tout sur leur passage. »}
\end{obstaclebox}

\begin{creaturebox}{première apparition -- les zèbres-crapauds}
\textit{« Leurs bulbes dorsaux gonflent et dégonflent en synchronisation, émettant une basse que vous ressentez dans vos os. Leurs yeux sont vides. »}
\end{creaturebox}

\begin{momentbox}{traversée -- dans la marée vivante}
\textit{« Des corps massifs vous frôlent de toutes parts. Les vibrations font trembler vos dents. Musc, herbe, quelque chose de primordial. »}
\end{momentbox}

\begin{obstaclebox}{découverte -- le dormeur}
\textit{« Recroquevillé comme un fœtus titanesque. Quinze mètres de chair pâle et veineuse, tentacules lovés, gueule fermée sur des dents comme des sabres. Elle dort et bloque le passage que vous pensiez prendre au milieu des herbes Vous ne voyez pas d'autre chemin. »}
\end{obstaclebox}

\begin{puzzlebox}{puzzle -- les fleurs-hélices}
\textit{« Autour de la créature, des plantes de deux mètres frémissent dans le vent. Leurs pétales tournent -- lentement d'abord, puis de plus en plus vite. L'une se déracine et s'élève dans les airs. Elle flotte. Rapidement, plusieurs d'entre elles s'agglutinent et elles semblent ne flotter que mieux. »}
\end{puzzlebox}

\begin{creaturebox}{première apparition -- le dormeur (si réveil)}
\textit{« Les veines s'illuminent d'abord. Un réseau de lumière bleue sous la surface. L'eau frémit. Des bulles. La silhouette se déplie -- lentement, majestueusement. Quinze mètres. Six tentacules qui se déploient comme les branches d'un arbre mort. La gueule s'ouvre. »}
\end{creaturebox}

\begin{creaturebox}{première apparition -- les fleurs-hélices}
\textit{« La plante vibre sous vos doigts. Ses pétales accélèrent. Elle tire vers le haut. »}
\end{creaturebox}

\begin{momentbox}{survol du dormeur}
\textit{« Le vent vous porte. Le lac défile en dessous, transparent. Chaque détail de la créature visible -- les cicatrices millénaires, les parasites dans ses plis. Ses yeux fermés semblent suivre votre ombre derrière leurs paupières. Ne faites pas de bruit. »}
\end{momentbox}

\begin{zonebox}{arrivée -- le vaisseau}
\textit{« À moitié enfoncé dans le sol, coque percée, envahi par des excroissances végétales. La sirène résonne faiblement. Des lumières clignotent derrière les hublots crasseux. La porte du sas est entrouverte. Quelque chose a forcé le passage. »}
\end{zonebox}

\begin{obstaclebox}{découverte -- l'intérieur du vaisseau}
\textit{« Lianes et racines dans les couloirs, pulsant faiblement. Éclairage d'urgence rouge. Écrans corrompus. Bips irréguliers du système de survie. »}
\end{obstaclebox}

\begin{creaturebox}{première apparition -- hollow}
\textit{« Dans la salle de cryogénisation, une silhouette immobile. Deux mètres. Lisse, noir et blanc en motifs organiques abstraits. Pas de visage -- une surface réfléchissante qui renvoie vos propres regards, déformés. Elle incline la tête. Une pression s'abat sur vos pensées. »}
\end{creaturebox}

\begin{creaturebox}{première apparition -- écho-écho (alternative)}
\textit{« Silence total dans les couloirs. Puis un claquement de dents. Rapide. Métallique. Silhouette humanoïde dans la pénombre -- pas de visage, une parabole organique. Sa peau ondule au moindre son. »}
\end{creaturebox}

\begin{momentbox}{possession par hollow}
\textit{(à lire au joueur possédé, en aparté)}
\textit{« Le silence. La paix. Tes pensées s'organisent avec une clarté parfaite. Hollow n'est pas ton ennemi. Hollow t'a libéré. Empêche les autres de lui faire du mal. »}
\end{momentbox}

\end{document}
